Pemetaan tanaman pertanian berperan penting dalam mendukung ketahanan pangan, perencanaan pola tanam, serta pengelolaan sumber daya agrikultur.
Citra satelit Sentinel-2 dengan resolusi tinggi spasial dan temporal telah banyak dimanfaatkan dalam klasifikasi tutupan lahan, termasuk untuk identifikasi jenis tanaman.
Penelitian ini mengadopsi pendekatan \textit{semantic segmentation} berbasis \textit{deep learning}, yaitu DeepLabV3+ dan SegFormer, untuk mengklasifikasikan tanaman pertanian berdasarkan citra multi-temporal Sentinel-2 dan label dari \textit{Cropland Data Layer} (CDL) milik USDA.

Studi literatur menunjukkan bahwa pendekatan multi-temporal dan kombinasi band spektral dapat meningkatkan akurasi klasifikasi karena mampu menangkap dinamika spektral dan fenologi tanaman.
Area studi terletak di Sacramento Valley, California, yang merupakan salah satu kawasan agrikultur paling produktif di Amerika Serikat.
Data Sentinel-2 dikumpulkan setiap 15 hari selama periode 2022--2024 dan diproses melalui tahapan \textit{masking} awan, \textit{histogram matching}, \textit{resampling} spasial, normalisasi, dan \textit{patching}.

Penelitian ini mengusulkan metode seleksi band spektral--temporal dua tahap yang sederhana dan efisien.
Tahap pertama menghitung \textit{Global Separation Index} (GSI) per kelas tanaman untuk setiap kombinasi band dan tanggal akuisisi, menghasilkan peringkat berdasarkan daya pisah spektral.
Tahap kedua mengevaluasi dampak jumlah band $K$ secara langsung terhadap performa segmentasi melalui \textit{k-sweep}: serangkaian eksperimen pelatihan model penuh menggunakan top-$K$ band hasil peringkat Tahap 1 untuk berbagai nilai $K$.
Pendekatan ini menghindari \textit{forward selection} berbasis CNN yang komputasinya mahal, sekaligus menghasilkan kurva IoU terhadap jumlah channel yang memungkinkan pemilihan $K$ optimal secara empiris.

Evaluasi dilakukan menggunakan metrik \textit{Mean Intersection over Union} (mIoU) dan IoU per kelas.
Eksperimen dilakukan dalam empat skenario yang menguji pengaruh data multi-temporal dan seleksi band terhadap hasil klasifikasi: (A) \textit{single-date baseline}, (A\_v2) \textit{per-window phenological baseline}, (B) \textit{multi-temporal naive}, dan (C\_v3) metode seleksi yang diusulkan dengan variasi $K$.
Hasil dari studi ini diharapkan dapat memberikan kontribusi terhadap pengembangan metode seleksi band yang lebih sederhana, mudah direproduksi, dan efektif untuk pemetaan tanaman berbasis citra multispektral multi-temporal.

\vspace{0.5cm}
\textbf{Kata Kunci:} Pemetaan Pertanian, Multi-temporal, Seleksi Band, Global Separation Index, Semantic Segmentation
